\documentclass[10pt]{article}
\usepackage[T1]{fontenc}
\usepackage[utf8]{inputenc}
\usepackage[french]{babel}
\usepackage{amsmath,amssymb,mathrsfs,mathpazo}
\usepackage{geometry}
\geometry{a4paper,margin=2cm}
\usepackage{graphicx,array,tabularx,tabularray,colortbl,multirow,multicol,booktabs,enumitem,pifont,ulem,fancybox,diagbox,cancel}
\usepackage{pstricks,pstricks-add,pst-plot,pst-tree,pst-node,pst-text,pst-func,pst-math,pst-bspline,pst-3dplot,pst-eucl,pst-solides3d}
\let\pscalculate\undefined
\usepackage{pst-bezier}
\usepackage{tikz,pgfplots,tkz-tab}
\usetikzlibrary{arrows,arrows.meta,patterns,calc,positioning,trees,decorations.markings,shapes.geometric,shapes.misc,shapes.arrows,intersections,decorations.pathmorphing,decorations.pathreplacing,backgrounds,shadows,babel}
\pgfplotsset{compat=1.18}
\usepgfplotslibrary{fillbetween}
\usepackage[np]{numprint}
\usepackage{listings,makecell,pas-tableur,twoopt}
\let\barre\undefined
\usepackage[pstricks]{tablvar}
\definecolor{darkgreen}{RGB}{0,100,0}
\providecommand{\red}{\color{red}}
\providecommand{\blue}{\color{blue}}
\providecommand{\green}{\color{green}}
\providecommand{\black}{\color{black}}
\providecommand{\white}{\color{white}}
\providecommand{\gray}{\color{gray}}
\providecommand{\cyan}{\color{cyan}}
\providecommand{\magenta}{\color{magenta}}
\providecommand{\yellow}{\color{yellow}}
\providecommand{\violet}{\color{violet}}
\providecommand{\orange}{\color{orange}}
\providecommand{\euro}{\texteuro}
\providecommand{\EUR}{\texteuro}
\providecommand{\ieme}{\textsuperscript{e}}
\providecommand{\degres}{\ensuremath{^\circ}}
\providecommand{\degree}{\ensuremath{^\circ}}
\providecommand{\vv}[1]{\overrightarrow{#1}}
\providecommand{\Coord}[1]{#1}
\providecommand{\decosix}{}
\providecommand{\starredbullet}{\ensuremath{\bullet}}
\pagestyle{empty}
\setlength{\parindent}{0pt}
\usepackage{adjustbox}\DeclareUnicodeCharacter{2236}{:}\DeclareUnicodeCharacter{2190}{\ensuremath{\leftarrow}}
\providecommand{\R}{\mathbb{R}}
\providecommand{\N}{\mathbb{N}}
\providecommand{\Z}{\mathbb{Z}}
\providecommand{\Q}{\mathbb{Q}}
\providecommand{\D}{\mathbb{D}}
\providecommand{\C}{\mathbb{C}}
\providecommand{\e}{{\mathrm{e}}}
\providecommand{\vect}[1]{\overrightarrow{#1}}
\providecommand{\vectt}[1]{\overrightarrow{\mathrm{#1}}}
\providecommand{\vvt}[1]{\overrightarrow{\mathrm{#1}}}
\providecommand{\np}[2][]{{#2~#1}}
\providecommand{\vv}[1]{{\overrightarrow{#1}}}
\renewcommand{\d}{{\mathrm{d}}}
\providecommand{\degres}{{^\circ}}
\providecommand{\degre}{{^\circ}}
\providecommand{\blue}{\color{blue}}
\providecommand{\red}{\color{red}}
\providecommand{\enskip}{\quad}
\providecommand{\llbracket}{\lbrack\!\lbrack}
\providecommand{\rrbracket}{\rbrack\!\rbrack}
\providecommand{\euro}{{\text{€}}}
\providecommand{\no}{{\text{n°}}}
\providecommand{\interval}[2]{{[#1;#2]}}
\providecommand{\intervFO}[2]{{[#1;#2[}}
\providecommand{\dotfill}{\ldots}

\providecommand{\ds}{\displaystyle}
\providecommand{\textding}[1]{#1}
\providecommand{\Oij}{\(\mathrm{O};\vec{\imath},\vec{\jmath}\)}
\providecommand{\Oijk}{\(\mathrm{O};\vec{\imath},\vec{\jmath},\vec{k}\)}
\providecommand{\Ouv}{\(\mathrm{O};\vec{u},\vec{v}\)}
\providecommand{\intervFF}[2]{\ensuremath{[#1;#2]}}
\renewcommand{\euro}{\ifmmode\text{\texteuro}\else\texteuro\fi}
\DeclareUnicodeCharacter{00B0}{\ifmmode^{\circ}\else\textdegree\fi}
\DeclareUnicodeCharacter{20AC}{\ifmmode\text{\texteuro}\else\texteuro\fi}
% Énoncés : Éducation nationale. Transcription et corrigés : APMEP, Denis Vergès et contributeurs. Découpage et classement : CoopMaths, Sébastien Lozano, Éric Elter et contributeurs. Adaptation de la présentation : Maths Explorer.
% CC BY-SA : https://coopmaths.fr/www/about/
% https://coopmaths.fr/alea/static/bac/2024/tex/bac_2024_06_sujet2_etrangers_3_cor.tex
\begin{document}
\noindent\textit{APMEP / CoopMaths. CC BY-SA. Présentation adaptée par Maths Explorer.}\par\medskip
\medskip


\begin{enumerate}
	\item

	On a : I$\left(\dfrac{1}{2}~;~ 0~;~0 \right)$  et J$\left(1~;~1~;~\dfrac{1}{2} \right)$ .

	\item

	On en déduit :

	$\vect{\text{EJ}} \begin{pmatrix}1-0  \\ 1- 0 \\ \dfrac{1}{2}-1     \end{pmatrix} =  \begin{pmatrix}1 \\ 1 \\ -\dfrac{1}{2}\end{pmatrix}$;
	\qquad
	$\vect{\text{FH}} \begin{pmatrix}0-1   \\ 1-0 \\ 1- 1   \end{pmatrix} =  \begin{pmatrix}-1 \\ 1 \\ 0\end{pmatrix}$;
	\quad	et \quad $\vect{\text{FI}} \begin{pmatrix}\dfrac{1}{2}-1  \\ 0-0  \\ 0-1   \end{pmatrix} =  \begin{pmatrix}-\dfrac{1}{2} \\ 0 \\ -1    \end{pmatrix}$.


	donc :

	$\vect{\text{EJ}} \cdot \vect{\text{FH}}= 1 \times (-1) + 1\times 1 -\dfrac{1}{2} \times 0 = 0$

	$\vect{\text{EJ}} \cdot \vect{\text{FI}}=  -\dfrac{1}{2} \times 1 + 1\times 0 -\dfrac{1}{2} \times (-1) = 0$

	Le vecteur $\vect{\text{EJ}}$ est donc orthogonal à deux vecteurs non colinéaires du plan (FHI) donc $\vect{\text{EJ}}$ est normal au plan (FHI).

	\item Le vecteur $\vect{\text{EJ}}\begin{pmatrix}1 \\ 1 \\ -\dfrac{1}{2}    \end{pmatrix}$ est normal au plan (FHI) donc le vecteur   $\vect{n}= -2\vect{\text{EJ}}\begin{pmatrix}-2 \\ -2 \\ 1    \end{pmatrix}$ est aussi un vecteur normal au plan (FHI).

Une équation cartésienne du plan (FHI)  est donc de la forme : \quad   $-2x-2y+z+d=0$ avec $d \in \R$.

De plus, le point F appartient au plan (FHI) donc ses coordonnées verifient l'équation du plan. On a donc :

$\begin{aligned}[t] \mathrm{F} \in (\mathrm{FHI}) &\iff -2\times x_\mathrm{F}  -2 \times y_\mathrm{F} + z_\mathrm{F} +d=0 \\
	&\iff -2\times 1  -2 \times 0 + 1 +d=0 \\
	&\iff -2+1+d=0 \\
	&\iff d=1\end{aligned}$

Une équation cartésienne du plan (FHI) est donc	$-2x-2y+ z+1=0$.
 \item Le vecteur $\vect{\text{EJ}} $ est un vecteur directeur de la droite (EJ) et E est un point de la droite (EJ).  Une représentation paramétrique de la droite (EJ) est donc :

$\begin{cases} x= 0 + 1\times t \\ y=0 + 1 \times t  \\ z=1 - \dfrac{1}{2} \times t  \end{cases} $ avec $t \in \R$ \quad c'est à dire \quad : $\begin{cases} x= t \\ y=t  \\ z=1 - \dfrac{1}{2} t  \end{cases} $ avec $t \in \R$

	\item
	\begin{enumerate}
		\item Soit K est le projeté orthogonal du point E sur le plan (FHI).

K est donc l'intersection du plan (FHI) et de la droite orthogonale au plan (FHI) passant par E, c'est à dire la droite (EJ).

Pour cela, on va considérer le point $M_t$ de paramètre $t$ sur la droite $(\mathrm{EJ})$.

$\begin{aligned} M_t \in (\mathrm{FHI}) &\iff -2x_{M_t} -2y_{M_t}+z_{M_t} +1 = 0\\
&\iff -2\times t -2  \times t  +  \left( 1- \dfrac{1}{2}t\right) +1=0 \\
&\iff -2t-2t+1-\dfrac{1}{2}t+1=0 \\
&\iff -\dfrac{9}{2}t = -2 \\
&\iff t= \dfrac{4}{9} \end{aligned}$

Le seul point de la droite étant aussi sur le plan est donc K, c'est le point de paramètre $\dfrac{4}{9}$ dans l'équation paramétrique de $(\mathrm{EJ})$.

		Ses coordonnées vérifient :

		$x_\mathrm{K}=\dfrac{4}{9}$, \quad $y_\mathrm{K}=\dfrac{4}{9}$ \quad et \quad $z_\mathrm{K}=1- \dfrac{1}{2} \times \dfrac{4}{9}=1-\dfrac{2}{9}=\dfrac{7}{9}$.

		Le projeté orthogonal  du point E sur le plan (FHI) est le point K de coordonnées K$\left( \dfrac{4}{9} ~;~  \dfrac{4}{9} ~;~ \dfrac{7}{9} \right)$.


		\item  Le volume $V$ d'une pyramide est donné par  $V=\dfrac{1}{3}\mathcal{B}h$  où  $\mathcal{B}$ est l'aire de la base et $h$ la hauteur correspondante.

		Prenons le triangle EFH comme base, la hauteur issue de  I est donc la droite (IL) et la distance de I au plan EFH est donc égale à la longueur IL.

		EFH est un triangle rectangle en E et son aire est égale à $\dfrac{ 1\times 1}{2}=\dfrac{1}{2}$

		L a pour coordonnées  L$\left(\dfrac{1}{2}~;~ 0~;~1 \right)$   donc

		LI$^2= \left( \dfrac{1}{2} -\dfrac{1}{2}   \right)^2 + (0-0)^2 + (1-0)^2= 1 $ donc LI$=1$.

		On a donc V= $\dfrac{1}{3} \times \dfrac{1}{2} \times 1 = \dfrac{1}{6}~\left(\text{cm}^3\right)$.


		\item  Calculons maintenant le volume en prenant le triangle FHI comme base. La hauteur sera donc la longueur EK.

		EK$^2=    \left( \dfrac{4}{9} -0   \right)^2 + \left( \dfrac{4}{9} -0  \right)^2 + \left( \dfrac{7}{9} -1  \right)^2  = \dfrac{16}{81}+ \dfrac{16}{81}+\dfrac{4}{81}  = \dfrac{36}{81}$

		donc EK$=\dfrac{6}{9}=\dfrac{2}{3}$

		On a donc :

		$\begin{aligned}[t] \dfrac{1}{6}=\dfrac{1}{3}\times \mathcal{A}_{\mathrm{FHI}}\times \dfrac{2}{3}
		&\iff \dfrac{1}{6} \times 3 \times\dfrac{3}{2} =  \mathcal{A}_{\mathrm{FHI}} \\
		&\iff \dfrac{1}{6} \times 3 \times\dfrac{3}{2} =  \mathcal{A}_{\mathrm{FHI}} \\
		&\iff  \mathcal{A}_{\mathrm{FHI}}=\dfrac{3}{4}\end{aligned}$


		L'aire du triangle FHI est	$ \dfrac{3}{4}~\left(\text{cm}^2\right)$.
	\end{enumerate}
\end{enumerate}

\bigskip

\end{document}
